% \iffalse meta-comment
%
% Copyright (C) 2026 by Your Name
%
% This file may be distributed and/or modified under the
% conditions of the LaTeX Project Public License, either version 1.3c
% of this license or (at your option) any later version.
%
% \fi
%
% \iffalse
%<*driver>
\ProvidesFile{zjureport.dtx}[2026/04/07 v2.1 ZJU Lab Report Template]
\documentclass{ctxdoc}
\usepackage{hypdoc}
\begin{document}
  \DocInput{zjureport.dtx}
\end{document}
%</driver>
% \fi
%
% \CheckSum{0}
%
% \CharacterTable
%  {Upper-case    \A\B\C\D\E\F\G\H\I\J\K\L\M\N\O\P\Q\R\S\T\U\V\W\X\Y\Z
%   Lower-case    \a\b\c\d\e\f\g\h\i\j\k\l\m\n\o\p\q\r\s\t\u\v\w\x\y\z
%   Digits        \0\1\2\3\4\5\6\7\8\9
%   Exclamation   \!     Double quote  \"     Hash (number) \#
%   Dollar        \$     Percent       \%     Ampersand     \&
%   Acute accent  \'     Left paren    \(     Right paren   \)
%   Asterisk      \* Plus          \+     Comma         \,
%   Minus         \-     Point         \.     Solidus       \/
%   Colon         \:     Semicolon     \;     Less than     \<
%   Equals        \=     Greater than  \>     Question mark \?
%   Commercial at \@     Left bracket  \[     Backslash     \\
%   Right bracket \]     Circumflex    \^     Underscore    \_
%   Grave accent  \`     Left brace    \{     Vertical bar  \|
%   Right brace   \}     Tilde         \~}
%
% \changes{v2.1}{2026/04/07}{恢复原始视觉布局，修正字体调用与变量定义}
%
% \GetFileInfo{zjureport.dtx}
%
% \title{浙江大学实验报告模板 \texttt{zjureport} \thanks{这份文档对应 \filename ~ \fileversion，最后修改日期 \filedate。}}
% \author{Your Name}
% \date{\filedate}
%
% \maketitle
%
% \section{简介}
% 本模板是为浙江大学本科生设计的实验报告 \LaTeX{} 文档类。基于 \texttt{ctexart} 编写。
%
% \section{使用方法}
% 提取文件后，请使用 \texttt{xelatex} 编译。
%
% \subsection{宏包选项}
% \DescribeMacro{codeformat} 
% 使用模板预设的 \texttt{listings} 代码格式。
%
% \subsection{用户命令}
% \DescribeMacro{\zjulogo} 
% 导入校徽图片。
%
% \DescribeMacro{\name} 
% 设置学生姓名。
%
% \DescribeMacro{\stuid} 
% 设置学生学号。
%
% \DescribeMacro{\makeheader} 
% 生成正文首页头部信息。
%
% \DescribeMacro{\makecover} 
% 生成独立封面。
%
% \StopEventually{}
%
% \section{代码实现}
%
%    \begin{macrocode}
%<*class>
\NeedsTeXFormat{LaTeX2e}
\ProvidesClass{zjureport}[2026/04/07 v2.1 ZJU Lab Report Template]

\newif\ifzju@codeformat
\zju@codeformatfalse
\DeclareOption{codeformat}{\zju@codeformattrue}
\DeclareOption*{\PassOptionsToClass{\CurrentOption}{ctexart}}
\ProcessOptions\relax

\LoadClass[fontset=fandol]{ctexart}

\RequirePackage{graphicx, xcolor, tikz}
\RequirePackage{geometry, titlesec, fancyhdr}
\RequirePackage{indentfirst, enumitem, zhnumber}

\linespread{1.2}
\geometry{top=1in,bottom=1in,left=1in,right=1in}

\setcounter{secnumdepth}{3}
\renewcommand\thesection{\zhnumber{\arabic{section}}、}
\renewcommand\thesubsection{\arabic{subsection}.}
\renewcommand\thesubsubsection{(\arabic{subsubsection})}
\setlist[enumerate,1]{label=(\arabic*)}
\setlist[enumerate,2]{label=\alph*.}
\titleformat*{\section}{\large\bfseries}
\titleformat*{\subsection}{\normalsize}
\titleformat*{\subsubsection}{\normalsize}

\ifzju@codeformat
  \RequirePackage{listings, cncolours}
  \colorlet{keyword}{松花绿}
  \colorlet{comment}{漆黑!50}
  \colorlet{emph1}{靛蓝}
  \colorlet{emph2}{琥珀}
  \lstset{
    frame=tb, framerule=0.5pt, numbers=none,
    breaklines=true, breakatwhitespace=true,
    keywordstyle = \bfseries\color{keyword},
    commentstyle = \itshape\color{comment},
    emphstyle    = [1]\itshape\color{emph1},
    emphstyle    = [2]\color{emph2},
    basicstyle   = {\small\ttfamily},
  }
\fi

\newcommand{\@major}{}
\newcommand{\@name}{}
\newcommand{\@stuid}{}
\newcommand{\@college}{}
\newcommand{\@lab}{}
\newcommand{\@course}{}
\newcommand{\@instructor}{}
\newcommand{\@grades}{}
\newcommand{\@expname}{}
\newcommand{\@exptype}{}
\newcommand{\@partner}{}

\newcommand\major[1]{\def\@major{#1}}
\newcommand\name[1]{\def\@name{#1}}
\newcommand\stuid[1]{\def\@stuid{#1}}
\newcommand\college[1]{\def\@college{#1}}
\renewcommand\date[1]{\def\@date{#1}}
\newcommand\lab[1]{\def\@lab{#1}}
\newcommand\course[1]{\def\@course{#1}}
\newcommand\instructor[1]{\def\@instructor{#1}}
\newcommand\grades[1]{\def\@grades{#1}}
\newcommand\expname[1]{\def\@expname{#1}}
\newcommand\exptype[1]{\def\@exptype{#1}}
\newcommand\partner[1]{\def\@partner{#1}}

\newcommand{\@zjulogo}[1]{%
  \IfFileExists{figures/浙江大学.pdf}{%
    \includegraphics[width=#1]{figures/浙江大学.pdf}%
  }{%
    \textcolor{black!15}{\rule{#1}{1.2cm}}% 根据传入的宽度渲染占位符
  }%
}
\newcommand{\zjulogo}[1]{%
  \renewcommand{\@zjulogo}[1]{\includegraphics[width=##1]{#1}}%
}
\pagestyle{fancy}
\lhead{实验名称：\@expname}
\chead{姓名：\@name}
\rhead{学号：\@stuid}

\def\zjuu[#1]#2{\underline{\hbox to #1{\hfil#2\hfil}}}

\newcommand{\makeheader}{
\begin{figure}[h]
  \begin{minipage}{0.4\linewidth}
    \centerline{\@zjulogo{\linewidth}}
  \end{minipage}%
  \begin{minipage}{0.3\linewidth}
    \centerline{\fontsize{36}{50}\selectfont\kaishu\addCJKfontfeatures{FakeBold=2.5} 实验报告}
  \end{minipage}%
  \begin{minipage}{0.3\linewidth}
    \raggedleft
    \begin{tabular*}{.8\linewidth}{ll}
      专业： & \zjuu[70pt]\@major \\
      姓名： & \zjuu[70pt]\@name  \\
      学号： & \zjuu[70pt]\@stuid \\
      日期： & \zjuu[70pt]\@date \\
      地点： & \zjuu[70pt]\@lab
    \end{tabular*}
  \end{minipage}
\end{figure}

\begin{table}[!htbp]
  \centering
  \begin{tabular*}{\linewidth}{llllll}
    课程名称： & \zjuu[150pt]\@course  & 
    指导老师： & \zjuu[50pt]\@instructor & 
    成绩：     & \zjuu[50pt]\@grades \\
    实验名称： & \zjuu[150pt]\@expname & 
    实验类型： & \zjuu[50pt]\@exptype    & 
    同组姓名： & \zjuu[50pt]\@partner
  \end{tabular*}
\end{table}
\thispagestyle{empty}
}

\newcommand{\makecover}{
  \begin{center}
    \begin{center}
      \@zjulogo{0.7\linewidth}
    \end{center}
    \vspace{1cm}
    \huge{\textbf{本科实验报告}}\\
    \vspace{0.5cm}
    \Large{\textbf{\@expname}}
    \vspace{3cm}
    \begin{center}
      \Large\linespread{2}\centering
      \renewcommand\arraystretch{1.6}
      \begin{tabular}{p{3cm}p{6cm}}
        \textbf{课程名称：} & \@course \\ \cline{2-2}
        \textbf{姓名：}     & \@name   \\ \cline{2-2}
        \textbf{学院：}     & \@college \\ \cline{2-2}
        \textbf{专业：}     & \@major  \\ \cline{2-2}
        \textbf{学号：}     & \@stuid  \\ \cline{2-2}
        \textbf{指导老师：} & \@instructor  \\ \cline{2-2}
      \end{tabular}
    \end{center}
    \vfill
    \Large{\@date}
  \end{center}
  \thispagestyle{empty}\newpage
}
%</class>
%    \end{macrocode}
%
% \iffalse
%<*example>
\documentclass[codeformat]{zjureport}
\usepackage{hyperref}
\hypersetup{colorlinks=true}

\major{信息工程}
\name{张扬}
\stuid{3171002333}
\college{信息与电子工程学院}
\date{\zhtoday}
\lab{东四-A-308}
\course{\LaTeX 模板使用}
\instructor{赵林平}
\grades{93}
\expname{使用实验报告模板}
\exptype{操作实验}
\partner{无}

\begin{document}
\makecover
\makeheader
\section{实验目的和要求}
测试楷体与非粗体封面内容效果。

\section{实验内容和步骤}

\subsection{实验内容}

\subsection{实验步骤}

\section{主要仪器设备}
计算机及配套软件。

\section{操作方法和实验步骤}

\section{实验数据记录和处理}

\begin{lstlisting}[language=C]
#include <stdio.h>
int main() {   
    int number;
   
    printf("Enter an integer: ");  
    
    // reads and stores input
    scanf("%d", &number);

    // displays output
    printf("You entered: %d", number);
    
    return 0;
}
\end{lstlisting}

\section{实验结果与分析}
\end{document}
%</example>
% \fi
%
% \Finale
\endinput